\documentclass[aspectratio=169, 11pt, t]{beamer}

% --- Theme and Colors ---
\usetheme{Boadilla}
\usecolortheme{default}
\setbeamertemplate{navigation symbols}{}
\setbeamertemplate{itemize items}[circle]
\setbeamertemplate{enumerate items}[default]

\definecolor{MyBlue}{RGB}{0, 51, 102}
\definecolor{BoxGrey}{RGB}{245, 245, 247}
\setbeamercolor{title}{fg=MyBlue}
\setbeamercolor{frametitle}{fg=MyBlue}
\setbeamercolor{structure}{fg=MyBlue}
\setbeamercolor{conclusionbox}{bg=BoxGrey, fg=black}

% --- Packages ---
\usepackage[utf8]{inputenc}
\usepackage[T1]{fontenc}
\usepackage{graphicx}
\usepackage{booktabs}
\usepackage{amsmath}
\usepackage{amsfonts}
\usepackage{amssymb}
\usepackage{caption}
\captionsetup{justification=raggedright,singlelinecheck=false,labelsep=colon,font=footnotesize,skip=-1pt}
\setbeamertemplate{caption}[numbered]
\usepackage{subcaption}
\usepackage[normalem]{ulem}
\usepackage{hyperref}
\hypersetup{
    colorlinks=true,
    linkcolor=blue,
    urlcolor=blue,
}

% --- Figure paths ---
\graphicspath{{./}{figures/}}

% --- Placeholder box for material to be added (optional arg: box height) ---
\newcommand{\placeholder}[2][0.52\textheight]{%
    \begin{center}
    \setlength{\fboxsep}{8pt}%
    \fcolorbox{gray}{BoxGrey}{%
        \begin{minipage}[c][#1][c]{0.95\textwidth}
            \centering\itshape\color{gray} #2
        \end{minipage}}%
    \end{center}}

% --- Presentation Info ---
\title[]{Jamaica Case Study}
\author{}
\institute{}
\date{\today}

% --- Layout ---
\setbeamertemplate{footline}[frame number]
\setbeamersize{text margin left=1cm,text margin right=1cm}

\begin{document}

% --- Title Slide ---
\begin{frame}[plain]
    \titlepage
\end{frame}

% --- Slide: reallocation of expenditure (FM panel 2) ---
\begin{frame}{Emerging Market and Developing Economies: Reallocation of Expenditure}
    \vspace{0.3em}
    \begin{columns}[T]
        \begin{column}{0.49\textwidth}
            \raggedright
            \captionof{figure}{Gains in output at selected horizons (percent deviation from steady state)}\label{fig:reallocationEM}
            {\centering\includegraphics[width=\textwidth]{reallocationEM_yd.png}\par}
        \end{column}
        \begin{column}{0.49\textwidth}
            \raggedright
            \captionof{figure}{Gains in output over time (percent deviation from steady state)}\label{fig:reallocationEMLines}
            {\centering\includegraphics[width=\textwidth]{reallocationEM_yd_lines.png}\par}
        \end{column}
    \end{columns}

    \vspace{1em}
    {\fontsize{5pt}{6pt}\selectfont
        Sources: IMF staff estimates. \\
        Notes: Output responses to a permanent increase, in the expenditure categories listed in the legends, of 1 percent of GDP in 2025, funded by a cut in public consumption. \\
    }
\end{frame}

% --- Slide: spending efficiency (FM panel 4) ---
\begin{frame}{Emerging Market and Developing Economies: Increase in Spending Efficiency}
    \vspace{0.3em}
    \begin{columns}[T]
        \begin{column}{0.49\textwidth}
            \raggedright
            \captionof{figure}{Additional gains in output in 2050 (percent deviation from steady state)}\label{fig:efficiencyEM}
            {\centering\includegraphics[width=\textwidth]{efficiencyEM_yd.png}\par}
        \end{column}
        \begin{column}{0.49\textwidth}
            \raggedright
            \captionof{figure}{Additional gains in output over time (percent deviation from steady state)}\label{fig:efficiencyEMLines}
            {\centering\includegraphics[width=\textwidth]{efficiencyEM_yd_lines.png}\par}
        \end{column}
    \end{columns}

    \vspace{1em}
    {\fontsize{5pt}{6pt}\selectfont
        Sources: IMF staff estimates. \\
        Notes: Additional gains in output from gradually closing gaps in spending efficiency, by 2040 and 2050, relative to the reallocation-only baseline. The left panel also shows gains from different initial efficiency levels, as described in its legend; the right panel shows the baseline case over time, with blue lines for infrastructure and purple lines for human capital investment. \\
    }
\end{frame}

% --- Slide: reallocation and efficiency combined, EMDEs versus Jamaica ---
\begin{frame}{Reallocation and Efficiency Combined: EMDEs and Jamaica}
    \vspace{0.3em}
    \begin{columns}[T]
        \begin{column}{0.49\textwidth}
            \raggedright
            \captionof{figure}{Emerging market and developing economies (percent deviation from steady state)}\label{fig:combinedEM}
            {\centering\includegraphics[width=\textwidth]{combinedEM_yd_lines.png}\par}
        \end{column}
        \begin{column}{0.49\textwidth}
            \raggedright
            \captionof{figure}{Jamaica (percent deviation from steady state)}\label{fig:combinedJAM}
            {\centering\includegraphics[width=\textwidth]{combinedJAM_yd_lines.png}\par}
        \end{column}
    \end{columns}

    \vspace{1em}
    {\fontsize{5pt}{6pt}\selectfont
        Sources: IMF staff estimates. \\
        Notes: Solid lines show output responses to a permanent reallocation of 1 percent of GDP in 2025 towards the expenditure category indicated by the color (blue: infrastructure investment; purple: human capital investment), funded by a cut in public consumption. Dashed lines show the overall effect when, in addition, the corresponding spending-efficiency gap is gradually closed by 2040. The Jamaica panel uses emerging-market parameters with Jamaica-specific efficiency gaps from IMF staff efficiency estimates (health 2023 and education 2024, averaged); the infrastructure gap uses the emerging-market average, as no Jamaica estimate is available. \\
    }
\end{frame}

\end{document}
